% Bagian: sets-functions-relations
% Bab: arithmetization
% Subbagian: cauchy

\documentclass[../../../include/open-logic-section]{subfiles}

\begin{document}

\olfileid[id]{sfr}{arith}{cauchy}

\olsection{Lampiran: Bilangan Real sebagai Barisan Cauchy}

Dalam \olref[cuts]{sec}, kita membangun bilangan real sebagai potongan
Dedekind. Dalam bagian ini, kita menjelaskan suatu konstruksi alternatif.
Konstruksi ini bertumpu pada definisi Cauchy tentang (apa yang kini kita sebut)
barisan Cauchy; tetapi penggunaan definisi ini untuk \emph{membangun} bilangan
real berasal dari para penulis lain pada abad ke-19, khususnya Weierstrass,
Heine, M\'{e}ray, dan Cantor. (Untuk sejarah yang menarik, lihat
\citeauthor{OConnorRobertson:RN} \citeyear{OConnorRobertson:RN}.)

Sebelum memasuki abad ke-19, ada baiknya kita mempertimbangkan Simon Stevin
(1548--1620). Singkatnya, Stevin menyadari bahwa kita dapat memikirkan setiap
bilangan real melalui ekspansi desimalnya. Maka, bahkan bilangan irasional
seperti $\sqrt{2}$ mempunyai ekspansi desimal yang baik, yang dimulai dengan:
\[
	1.41421356237\ldots
\]
Sangat mudah memodelkan ekspansi desimal dalam teori himpunan: cukup pandang
ekspansi itu sebagai fungsi $d \colon \Nat \to \Nat$, dengan $d(n)$ yang
menyatakan digit pada tempat desimal ke-$n$ yang kita perhatikan. Kemudian kita memerlukan sedikit
penyesuaian untuk menangani bagian bilangan real sebelum tanda desimal (di
sini, hanya $1$). Kita juga memerlukan penyesuaian lebih lanjut (suatu relasi
ekuivalensi) untuk menjamin bahwa, misalnya, $0.999\ldots = 1$. Namun, tidak
sulit memberikan konstruksi bilangan real yang sepenuhnya ketat secara
matematis, dengan cara Stevin, di dalam teori himpunan.

Stevin bukan fokus kita. (Untuk pembahasan lebih lanjut mengenai Stevin, lihat
\citealt{KatzKatz2012}.) Namun, berikut suatu gagasan yang berkaitan erat.
Alih-alih memperlakukan ekspansi desimal $\sqrt{2}$ secara langsung, kita
dapat mempertimbangkan suatu \emph{barisan} hampiran rasional yang semakin
akurat terhadap $\sqrt{2}$, dengan memperhatikan ekspansi yang semakin teliti:
\[
1, 1.4, 1.414, 1.4142, 1.41421,\ldots
\]
Gagasan bahwa bilangan real dapat dipandang melalui ``hampiran yang semakin
baik'' memberi kita landasan bagi rangkaian wawasan lainnya (serupa dengan
pengamatan yang kita gunakan ketika membangun $\Rat$ dari $\Int$, atau $\Int$
dari $\Nat$). Wawasan dasarnya adalah:
\begin{enumerate}
	\item Setiap bilangan real dapat ditulis sebagai ekspansi desimal
	(yang mungkin tak berhingga).
	\item Informasi yang dikodekan oleh ekspansi desimal (yang mungkin tak
	berhingga) dapat pula dikodekan oleh suatu barisan bilangan rasional.
	\item Suatu barisan bilangan rasional dapat dipandang sebagai fungsi
	dari $\Nat$ ke $\Rat$; cukup tetapkan $f(n)$ sebagai bilangan rasional
	ke-$n$ dalam barisan tersebut.
\end{enumerate}
Tentu saja, tidak \emph{setiap} fungsi dari $\Nat$ ke $\Rat$ akan memberi kita
suatu bilangan real. Sebagai contoh, perhatikan fungsi berikut:
\[
	f(n) =\begin{cases}
			1 & \text{jika }n\text{ ganjil}\\
			0 &\text{jika }n\text{ genap}
		\end{cases}
\]
Pada dasarnya, kekhawatirannya ialah bahwa barisan $0,1,0,1,0,1,0,\ldots$
tampaknya tidak ``mendekati'' bilangan real mana pun. Jadi, untuk memastikan
bahwa kita mempertimbangkan barisan yang memang mendekati suatu bilangan real,
kita perlu membatasi perhatian pada barisan yang mempunyai suatu
\emph{limit}.

Kita telah menjumpai gagasan limit dalam \olref[his][set][limits]{sec}. Namun,
kita tidak dapat menggunakan definisi yang \emph{persis} sama seperti yang
kita gunakan di sana. Ekspresi ``$(\forall \epsilon>0)$'' di sana secara
tersirat melibatkan kuantifikasi atas bilangan real; dan kita sedang
mempertimbangkan limit fungsi pada bilangan real; jadi, menggunakan definisi
tersebut berarti memanfaatkan bilangan real; padahal bilangan real itulah yang
sedang kita upayakan untuk \emph{bangun}. Untungnya, kita dapat bekerja dengan
gagasan limit yang berkaitan erat.
\begin{defn}\ollabel{def:CauchySequence}
  Suatu fungsi $f: \Nat \to \Rat$ merupakan \emph{barisan Cauchy} jika dan hanya jika untuk
  setiap $\epsilon \in \Rat$ positif berlaku $(\exists \ell \in
  \Nat)(\forall m, n > \ell)|f(m) - f(n)| < \epsilon$.
\end{defn}

Gagasan umum mengenai limit tetap sama seperti sebelumnya: jika Anda
menginginkan tingkat ketelitian tertentu (yang diukur oleh~$\epsilon$), ada
suatu ``wilayah'' tempat mencarinya (setiap masukan yang lebih besar
daripada~$\ell$). Mudah dilihat bahwa barisan kita $1$, $1.4$, $1.414$,
$1.4142$, $1.41421$\ldots mempunyai suatu limit: jika Anda ingin menghampiri
$\sqrt{2}$ dengan galat kurang dari $\nicefrac{1}{10^{n}}$, cukup perhatikan
sembarang suku setelah suku ke-$n$.

Gagasan yang jelas, dengan demikian, ialah mengatakan bahwa suatu bilangan
real \emph{adalah} sembarang barisan Cauchy. Namun, seperti dalam konstruksi
$\Int$ dan $\Rat$, gagasan ini terlalu naif: untuk setiap bilangan real
tertentu, ada beberapa barisan Cauchy berbeda yang menunjukkan bilangan real
tersebut. Cara sederhana untuk melihat hal ini adalah sebagai berikut. Jika
diberikan suatu barisan Cauchy~$f$, definisikan $g$ sebagai fungsi yang sama
persis dengan~$f$, kecuali bahwa $g(0)\neq f(0)$. Karena kedua barisan itu
setelah suku pertama, selisih keduanya mendekati nol menurut kriteria yang
berkaitan dengan \olref{def:CauchySequence}; karena itu, keduanya seharusnya
dipandang ``mendefinisikan'' bilangan real yang sama. Jadi, sebenarnya kita
harus memandang barisan-barisan Cauchy tersebut sebagai bilangan real yang
sama.

Oleh karena itu, sekali lagi kita perlu mendefinisikan suatu relasi ekuivalensi
pada barisan-barisan Cauchy, dan mengidentifikasi bilangan real dengan
kelas-kelas ekuivalensi. Pertama-tama, kita memerlukan gagasan tentang suatu fungsi yang
mendekati $0$ pada limitnya. Untuk setiap fungsi $h : \Nat \to \Rat$,
katakan bahwa \emph{$h$ mendekati $0$} jika dan hanya jika untuk setiap
$\epsilon \in \Rat$ positif berlaku $(\exists \ell \in \Nat)(\forall n >
\ell)|h(n)| < \epsilon$.\footnote{Bandingkan ini dengan definisi
$\lim_{x \mathord{\rightarrow}\infty}f(x) = 0$ dalam
\olref[his][set][limits]{sec}.} Selanjutnya, jika $f$ dan $g$ adalah fungsi
$\Nat \to \Rat$, tetapkan $(f-g)(n) = f(n) - g(n)$. Sekarang definisikan:
\[
	f \Realequiv g \text{ jika dan hanya jika $(f-g)$ mendekati $0$}.
\]
Kita perlu memeriksa bahwa $\Realequiv$ merupakan relasi ekuivalensi; dan
memang demikian. Kemudian, jika kita mau, kita dapat mendefinisikan bilangan
real sebagai kelas-kelas ekuivalensi, di bawah $\Realequiv$, dari semua
barisan Cauchy dari $\Nat \to \Rat$.

\begin{prob}
Misalkan $f(n) = 0$ untuk setiap $n$. Misalkan $g(n) = \frac{1}{(n+1)^2}$.
Tunjukkan bahwa keduanya merupakan barisan Cauchy, dan bahwa limit kedua
fungsi tersebut memang $0$, sehingga juga $f \sim_\Real g$.
\end{prob}

Setelah melakukan ini, seperti biasa kita menulis
$\equivrep{f}{\Realequiv}$ untuk kelas ekuivalensi yang memuat $f$ sebagai
!!a{element}. Namun, agar tetap mudah dibaca, selanjutnya kita akan
menghilangkan subskrip dan hanya menulis $\equivrep{f}{}$. Kita juga
menetapkan bahwa, untuk setiap $q \in \Rat$, berlaku $q_{\Real} =
\equivrep{c_{q}}{}$, dengan $c_{q}$ sebagai fungsi konstan $c_q(n) = q$ untuk
semua $n \in \Nat$. Kemudian kita mendefinisikan relasi dan operasi dasar pada
bilangan real, misalnya:
\begin{align*}
	\equivrep{f}{} + \equivrep{g}{} &= 	\equivrep{(f + g)}{} \\
	\equivrep{f}{} \times 	\equivrep{g}{} &= \equivrep{(f \times g)}{} 
\end{align*}
dengan $(f + g)(n) = f(n) + g(n)$ dan $(f \times g)(n) = f(n) \times
g(n)$. Tentu saja, kita juga perlu memeriksa bahwa masing-masing $(f + g)$,
$(f-g)$, dan $(f\times g)$ merupakan barisan Cauchy jika $f$ dan $g$ demikian;
memang benar, dan hal ini kita serahkan kepada Anda.

Terakhir, kita mendefinisikan suatu gagasan urutan. Katakan bahwa
$\equivrep{f}{}$ \emph{positif} jika dan hanya jika
$\equivrep{f}{}\neq 0_\Real$ dan $(\exists \ell \in \Nat)(\forall n >
\ell)0 < f(n)$ keduanya berlaku. Kemudian katakan bahwa $\equivrep{f}{} <
\equivrep{g}{}$ jika dan hanya jika $\equivrep{(g - f)}{}$ positif. Kita harus
memeriksa bahwa definisi ini terdefinisi dengan baik (yakni, tidak bergantung
pada pemilihan fungsi ``wakil'' dari kelas ekuivalensi).
%\begin{align*}
%	\equivrep{f}{} \leq \equivrep{g}{} \text{ jika dan hanya jika }&\equivrep{f}{} = \equivrep{g}{} \text{ atau }\\
%	&(\exists \ell \in \Nat)(\forall n \in \Nat)(\ell < n \lif f(n) \leq g(n))
%\end{align*}
Setelah melakukan ini, cukup mudah menunjukkan bahwa semuanya menghasilkan
sifat-sifat aljabar yang tepat; yakni:
\begin{thm}\ollabel{thm:cauchyorderedfield}
Barisan-barisan Cauchy membentuk suatu lapangan terurut.
\end{thm}

\begin{proof}
Latihan.
\end{proof}

\begin{prob}
Buktikan bahwa barisan-barisan Cauchy membentuk suatu lapangan terurut.
\end{prob}

Membuktikan bahwa bilangan real yang dibangun demikian memiliki Sifat
Kelengkapan lebih sulit, sehingga kita akan memberikan pembuktiannya.

\begin{thm}
Setiap himpunan takkosong barisan Cauchy yang mempunyai suatu batas atas
mempunyai suatu batas atas terkecil.
\end{thm}

\begin{proof}[Sketsa pembuktian] Misalkan $S$ adalah sembarang himpunan
takkosong barisan Cauchy yang mempunyai suatu batas atas. Maka terdapat suatu
$p \in \Rat$ sedemikian sehingga $p_{\Real}$ merupakan batas atas bagi $S$.
Ambil $r \in S$; maka terdapat suatu $q \in \Rat$ sedemikian sehingga
$q_{\Real} < \equivrep{r}{}$. Jadi, jika batas atas terkecil bagi $S$ ada, batas itu berada
di antara $q_\Real$ dan $p_\Real$ (inklusif).

Kita akan semakin mendekati b.a.t. tersebut secara bersamaan dari bawah dan
atas. Secara khusus, kita mendefinisikan dua fungsi, $f, g \colon \Nat \to
\Rat$, dengan tujuan agar $f$ semakin mendekati b.a.t.\ dari atas, dan $g$
semakin mendekatinya dari bawah. Kita mulai dengan mendefinisikan:
\begin{align*}
	f(0) &= p \\
	g(0) &= q
\end{align*}
Kemudian, dengan $a_n = \frac{f(n) + g(n)}{2}$, tetapkan:\footnote{Ini
merupakan definisi rekursif. Namun, kita \emph{belum} memberikan alasan apa
pun untuk menganggap bahwa definisi rekursif diperbolehkan.}
\begin{align*}
	f(n+1) &=
	\begin{cases}
		a_n &\text{jika }(\forall h \in S)\equivrep{h}{} \leq (a_n)_\Real\\
		f(n)&\text{jika tidak}
	\end{cases}\\
	g(n+1) &=
	\begin{cases}
		a_n &\text{jika }(\exists h \in S)\equivrep{h}{} \geq (a_n)_\Real\\
	 	g(n) &\text{jika tidak}
	\end{cases}
\end{align*}
Baik $f$ maupun $g$ merupakan barisan Cauchy. (Hal ini dapat diperiksa dengan
cukup mudah, tetapi kita serahkan sebagai latihan.) Perhatikan bahwa fungsi
$(f-g)$ mendekati $0$, karena selisih antara $f$ dan $g$ menjadi paling banyak setengah pada
setiap langkah. Oleh karena itu, $\equivrep{f}{} = \equivrep{g}{}$.

Pertama-tama kita menunjukkan bahwa $\equivrep{f}{}$ merupakan suatu batas atas bagi $S$, yakni\ $(\forall h \in S)\equivrep{h}{} \leq \equivrep{f}{}$.
(Selama pembuktian, kita akan menggunakan \olref{thm:cauchyorderedfield}.) Ambil $h \in S$ dan
andaikan, untuk memperoleh kontradiksi, bahwa $\equivrep{f}{} < \equivrep{h}{}$, sehingga
$0_\Real < \equivrep{(h-f)}{}$. Karena $f$ adalah barisan Cauchy yang menurun
secara monoton, terdapat suatu $n \in \Nat$ sedemikian sehingga
$\equivrep{(c_{f(n)} - f)}{} < \equivrep{(h-f)}{}$. Jadi:
\[
	(f(n))_\Real = \equivrep{c_{f(n)}}{} < \equivrep{f}{} + \equivrep{(h-f)}{} = \equivrep{h}{},
\]
yang bertentangan dengan fakta bahwa, menurut konstruksi, $\equivrep{h}{} \leq (f(n))_\Real$.

Selanjutnya kita menunjukkan bahwa $\equivrep{f}{} = \equivrep{g}{}$ adalah batas atas \emph{terkecil} bagi $S$. Jadi, misalkan $j$ adalah sembarang barisan Cauchy dan andaikan $\equivrep{j}{} < \equivrep{g}{}$. Dengan bernalar seperti di atas (menggunakan fakta bahwa $g$ \emph{meningkat}), terdapat $n \in \Nat$ sedemikian sehingga $\equivrep{j}{} < (g(n))_\Real$. Namun, menurut konstruksi terdapat $h \in S$ sedemikian sehingga $(g(n))_\Real \leq \equivrep{h}{}$, sehingga $\equivrep{j}{} < \equivrep{h}{}$ dan oleh karena itu $\equivrep{j}{}$ bukan batas atas bagi $S$.
\end{proof}

\end{document}
